\documentclass[12pt]{article}

\usepackage{amsmath}
\usepackage{listings}

\title{Assignment 1 Workout}
\author{Thomas Vroom - i6291496}
\date{September 2022}

\begin{document}
\maketitle
\pagenumbering{gobble}

\section*{Question 1}
Let $X$ be the number of people who show up to the movie.
Assuming the movie theater has sold 100 tickets and the 
probability of a single person showing up is 50\%, 
we can say:

\[ X \sim \text{Bin} (100, 0.5) \]

\smallskip
\noindent
To calculate the probability of $k$ people showing up to
the movie, we can use the following formula:

\[ P(X = k) = \binom{100}{k} {0.5}^k ({0.5}^{100-k}) \]

\smallskip
\noindent
Which can be simplified to:

\[ P(X = k) = \binom{100}{k} {0.5}^{100} \]

\smallskip
\noindent
Using this, we can calculate the probability of more than
50 people showing up to be:

\[ P(X > 50) = P(X = 51) + P(X = 52) + ... + P(X = 100) \]

\bigskip
\noindent
To calculate this we can use the following R code:

\begin{lstlisting}
    sum = 0;
    for (k in 51:100) {
        sum = sum + choose(100, k) * (0.5)^100;
    }
    print(sum);
\end{lstlisting}

\smallskip
\noindent
This gives us a final probability of $0.4602054$.

\section*{Question 2}

\textbf{1. Two random variables that are 
independent and identically distributed}

\smallskip
\noindent
X: Throw a dice.
Y: Throw another dice.

\bigskip
\noindent
\textbf{2. Two random variables that are 
independent and not identically distributed}

\smallskip
\noindent
X: Throw a 6-sided-dice.
Y: Throw an 8-sided-dice.

\bigskip
\noindent
\textbf{3. Two random variables that are 
dependent and identically distributed}

\smallskip
\noindent
X: Throw a dice.
Y: The result of X.

\bigskip
\noindent
\textbf{4. Two random variables that are 
dependent and not identically distributed}

\smallskip
\noindent
X: Throw a 6-sided-dice.
Y: If the result of X is 3 or lower, throw another
6-sided-dice. Else throw an 8-sided-dice.

\end{document}
